\section{Number of Provinces}
\label{sec:graph-number-of-provinces}

\problemstatement{Given an $n\times n$ adjacency matrix
\textit{isConnected}, where $\textit{isConnected}[i][j]=1$ means cities
$i$ and $j$ are directly connected, a \emph{province} is a maximal group
of directly or indirectly connected cities. Return the number of
provinces.}

\begin{patternbox}
\emph{``Count the connected components / groups''} is the plainest use
of \textbf{Disjoint Set}: union every directly connected pair, and the
number of provinces is simply the number of \textbf{distinct roots}
left. (A DFS/BFS component count also works; DSU is shown here because
this chapter is about wielding it.)
\end{patternbox}

\subsection*{From matrix to unions}

Scan the upper triangle of the matrix. Every $1$ at $(i,j)$ says ``$i$
and $j$ belong to the same province,'' which is exactly a
$\textit{union}(i,j)$. After processing all pairs, cities in the same
province share one root, so counting how many cities are their own root
(\,$\textit{find}(i)=i$\,) counts the provinces.

\begin{ideabox}[Components = Number of Distinct Roots]
Disjoint Set never explicitly stores group membership lists; it stores a
forest. After all unions, each tree is one component, and each tree has
exactly one root. So iterating once and counting self-parents yields the
component count in $O(n\,\alpha(n))$.
\end{ideabox}

\subsection*{Algorithm}

\begin{enumerate}
  \item Create a Disjoint Set over $n$ cities.
  \item For every pair $i<j$ with $\textit{isConnected}[i][j]=1$, call
        $\textit{unite}(i,j)$.
  \item Count the indices $i$ with $\textit{find}(i)=i$; that is the
        number of provinces.
\end{enumerate}

\subsection*{C++ implementation}

\begin{lstlisting}
#include <bits/stdc++.h>
using namespace std;

class DisjointSet {
    vector<int> parent, size;
public:
    DisjointSet(int n) : parent(n), size(n, 1) {
        for (int i = 0; i < n; i++) parent[i] = i;
    }
    int find(int x) {
        if (parent[x] == x) return x;
        return parent[x] = find(parent[x]);        // path compression
    }
    void unite(int x, int y) {
        int rx = find(x), ry = find(y);
        if (rx == ry) return;
        if (size[rx] < size[ry]) swap(rx, ry);
        parent[ry] = rx; size[rx] += size[ry];     // smaller under larger
    }
};

int numProvinces(vector<vector<int>>& isConnected) {
    int n = isConnected.size();
    DisjointSet ds(n);
    for (int i = 0; i < n; i++)
        for (int j = i + 1; j < n; j++)
            if (isConnected[i][j] == 1) ds.unite(i, j);

    int provinces = 0;
    for (int i = 0; i < n; i++) if (ds.find(i) == i) provinces++;
    return provinces;
}

int main() {
    vector<vector<int>> isConnected = {
        {1, 1, 0},
        {1, 1, 0},
        {0, 0, 1}
    };
    cout << numProvinces(isConnected) << endl;     // 2
    return 0;
}
\end{lstlisting}

\subsection*{Dry run}

\dryrun{The $3\times3$ matrix above.}

\begin{itemize}
  \item Pair $(0,1)$ is $1$: $\textit{unite}(0,1)$ -- $1$'s root becomes
        $0$.
  \item Pair $(0,2)$ and $(1,2)$ are $0$: no union.
  \item Roots: $\textit{find}(0)=0$, $\textit{find}(1)=0$,
        $\textit{find}(2)=2$. Self-parents are $0$ and $2$ -- two
        provinces.
\end{itemize}

\subsection*{Complexity}

\begin{complexitybox}
\textbf{Time Complexity.} $O(n^2\,\alpha(n))$ -- scanning the matrix
dominates, with near-constant union/find. \textbf{Space Complexity.}
$O(n)$ for the Disjoint Set arrays.
\end{complexitybox}

\begin{takeawaybox}
Counting provinces is Disjoint Set at its most direct: union every
connected pair, then count distinct roots. Recognize ``how many separate
groups?'' as a component-count question and DSU (or a DFS/BFS flood)
answers it immediately.
\end{takeawaybox}
